\documentclass[11pt]{article}

\usepackage[T1]{fontenc}

\usepackage[utf8]{inputenc}

\usepackage[margin=1in]{geometry}

\usepackage{graphicx}

\usepackage{amsmath}

\usepackage{amssymb}

\usepackage{booktabs}

\usepackage{array}

\usepackage{tabularx}

\usepackage{caption}

\usepackage{float}

\usepackage{microtype}

\usepackage{mathptmx}

\usepackage{natbib}

\usepackage[hidelinks]{hyperref}

\captionsetup{font=small,labelfont=bf}

\setlength{\parindent}{1.25em}

\setlength{\parskip}{0pt}

\title{Feasibility assessment of estimating e-cigarette prevalence in women under 50 and its association with epilepsy using the provided data hints}

\author{}

\date{}

\begin{document}

\maketitle

\begin{center}\small\textit{Research Memo}\end{center}

\begin{abstract}

This memo assessed whether the datasets provided in this run could support an estimate of e-cigarette use prevalence among women younger than 50 years and a test of association with epilepsy. The screening process focused on the BRFSS 2015 GitHub mirror named in the hints and on the other supplied dataset hints. The accessible BRFSS mirror was found to be a diabetes-focused derived dataset rather than a full BRFSS file with the needed e-cigarette and epilepsy variables \cite{ref1}. The other supplied dataset hints were not relevant to the research question or were outside the epidemiologic scope of population tobacco surveillance \cite{ref2}. Because no provided source could be verified as containing both target variables in analyzable form, no original empirical estimate was produced in this run.

\end{abstract}

\section{Introduction}
The research question was whether the datasets provided in this run could be used to estimate the prevalence of e-cigarette use among women under age 50 and examine whether e-cigarette use is associated with epilepsy. Answering that question required a population sample with both exposure and outcome variables, plus a way to restrict analyses to the target subgroup. The work performed here was a feasibility assessment of the available data hints rather than an epidemiologic analysis of observed associations. The goal was to determine whether any supplied source could honestly support the requested estimate and association test without introducing unsupported assumptions.

\section{Methods}
I converted the user request into a cross-sectional epidemiology question and screened the supplied dataset hints for relevance and variable availability. I searched for the BRFSS 2015 GitHub mirror named in the hints and reviewed accessible repository metadata and dataset descriptions to determine whether the mirrored file exposed variables for e-cigarette use and epilepsy. I also assessed the remaining hints, including MAST, DANDI, and NCI GDC, for topical suitability to a population-level question about e-cigarette prevalence and epilepsy in women under 50. I recorded the screening process in a run log and in a source suitability table \ref{tab:artifact-1} \ref{tab:artifact-2}. I did not compute prevalence, effect estimates, or models because no provided source could be verified as containing both target variables in analyzable form within the available environment.

\section{Results}
The source screening found that the accessible BRFSS 2015 mirror was not a full BRFSS surveillance file for this question. The repository associated with the mirror described a diabetes-focused derived dataset based on CDC BRFSS 2015, with variables oriented toward diabetes prediction rather than e-cigarette use and epilepsy \cite{ref1}. The search summary likewise indicated that the BRFSS mirror was insufficient and that the other provided datasets were not relevant to the research question \cite{ref2}. The artifact table documenting the run log shows the key empirical pattern of this run: the available materials did not expose a dataset suitable for honest estimation, and internet access constraints prevented direct downloading and inspection of remote files \ref{tab:artifact-1}. The source suitability table shows the same pattern across hints: one dataset was insufficient for the target question, while the others were out of scope \ref{tab:artifact-2}. As a result, no original prevalence estimate or association analysis was produced.

\begin{table}[tbp]
\centering
\caption{Run log documenting dataset screening and the reason no empirical estimate was produced.}
\label{tab:artifact-1}
\begin{tabular}{>{\raggedright\arraybackslash}p{0.15\linewidth}}
\toprule
Searched provided dataset hints and web sources. \\
\midrule
BRFSS 2015 diabetes mirror identified \\
No provided dataset hint contained a direct source suitable for estimating e-cig \\
Python runtime has no outbound internet \\
Therefore no original empirical estimate could be computed honestly in this run. \\
\bottomrule
\end{tabular}
\end{table}

\begin{table}[tbp]
\centering
\caption{Table summarizing suitability of each provided dataset hint for the user's research question.}
\label{tab:artifact-2}
\begin{tabular}{>{\raggedright\arraybackslash}p{0.15\linewidth}>{\raggedright\arraybackslash}p{0.15\linewidth}>{\raggedright\arraybackslash}p{0.15\linewidth}}
\toprule
dataset\_id & suitability & reason \\
\midrule
brfss-2015-github-mirror & insufficient & accessible mirror appears diabetes-focused; no e-cigarette or epilepsy variables \\
mast-observations & not relevant & astronomy archive unrelated to public-health question \\
dandi-api & not relevant & neuroscience archive may include epilepsy data but not population ecig prevalenc \\
nci-gdc-api & not relevant & oncology metadata unrelated to ecig prevalence and epilepsy in general populatio \\
\bottomrule
\end{tabular}
\end{table}

\section{Discussion}
The main conclusion from this run is not about the epidemiology of e-cigarette use and epilepsy, but about data feasibility. The supplied materials did not include a verified population dataset with both an e-cigarette-use variable and an epilepsy variable for women younger than 50. In that setting, attempting to report numeric prevalence or association results would have been misleading. The screening also showed that the most plausible hint, the BRFSS 2015 mirror, resolved to a diabetes-oriented derivative rather than the full surveillance resource needed for the research question \cite{ref1}. This makes the memo useful as a documentation of why the analysis could not proceed, and it identifies the kind of source that would be needed in a future run: a population-based dataset with direct measurement of tobacco use and epilepsy, plus demographic variables for subgroup restriction.

\section{Limitations}
No provided dataset hint could be verified in this run as containing both e-cigarette-use and epilepsy variables for women under age 50 in a population sample. The BRFSS 2015 GitHub mirror available through the hint appeared to be a diabetes-focused derived dataset, not the full BRFSS instrument needed for this question. The other supplied datasets were outside the scope of the epidemiologic question: MAST is astronomy, DANDI is neuroscience archive metadata rather than population tobacco surveillance, and NCI GDC is oncology-focused. The Python execution environment in this run could not directly download internet-hosted files, so candidate remote CSV contents could not be materialized and inspected locally after discovery. Because meaningful original estimation was not possible with the available materials, returning numeric prevalence or association results would have been misleading.

\bibliographystyle{plainnat}
\bibliography{references}

\end{document}
